
\documentclass[11pt,letterpaper,oneside]{report}
\usepackage[T1]{fontenc}
\usepackage[utf8]{inputenc}
\usepackage{lmodern}
\usepackage{inconsolata}
\usepackage{geometry}
\geometry{margin=0.89in,headheight=16pt,headsep=19pt,footskip=26pt}
\usepackage{xcolor}
\definecolor{navy}{HTML}{173758}
\definecolor{teal}{HTML}{157B89}
\definecolor{pale}{HTML}{E9F0F3}
\definecolor{codegray}{HTML}{F5F7F9}
\definecolor{muted}{HTML}{667785}
\usepackage{booktabs,longtable,tabularx,array,calc}
\usepackage{amsmath,amssymb}
\usepackage{graphicx}
\usepackage{microtype}
\usepackage{enumitem}
\usepackage{fancyhdr}
\usepackage{titlesec}
\usepackage{listings}
\usepackage{upquote}
\usepackage{xurl}
\usepackage{hyperref}
\hypersetup{colorlinks=true,linkcolor=navy,urlcolor=teal,bookmarksopen=true,pdfauthor={Software Build Engineering Documentation Suite}}
\urlstyle{same}
\setcounter{tocdepth}{2}
\setcounter{secnumdepth}{2}
\setlength{\parindent}{0pt}
\setlength{\parskip}{5pt plus 1pt}
\setlength{\emergencystretch}{3em}
\setlist[itemize]{leftmargin=*,topsep=4pt,itemsep=2pt}
\setlist[enumerate]{leftmargin=*,topsep=4pt,itemsep=2pt}
\titleformat{\chapter}[display]{\normalfont\bfseries\color{navy}}{\fontsize{11}{12}\selectfont CHAPTER \thechapter}{7pt}{\fontsize{23}{28}\selectfont}
\titlespacing*{\chapter}{0pt}{-6pt}{19pt}
\titleformat{\section}{\Large\bfseries\color{navy}}{\thesection}{0.65em}{}
\titleformat{\subsection}{\large\bfseries\color{teal}}{\thesubsection}{0.65em}{}
\pagestyle{fancy}
\fancyhf{}
\fancyhead[L]{\small\textcolor{muted}{PROFESSIONAL CMAKE ENGINEERING}}
\fancyhead[R]{\small\textcolor{muted}{\nouppercase{\leftmark}}}
\fancyfoot[C]{\small\textcolor{muted}{\thepage}}
\renewcommand{\headrulewidth}{0.3pt}
\fancypagestyle{plain}{\fancyhf{}\fancyfoot[C]{\small\thepage}\renewcommand{\headrulewidth}{0pt}}
\lstdefinelanguage{cmake}{sensitive=false,morekeywords={cmake_minimum_required,project,add_library,add_executable,target_compile_features,target_include_directories,target_link_libraries,target_sources,target_compile_options,target_compile_definitions,target_link_options,add_subdirectory,include,if,else,elseif,endif,foreach,endforeach,function,endfunction,macro,endmacro,set,option,find_package,find_program,add_custom_command,add_custom_target,install,configure_file,configure_package_config_file,write_basic_package_version_file,set_target_properties,set_tests_properties,add_test,message,FetchContent_Declare,FetchContent_MakeAvailable,return,cmake_parse_arguments},morecomment=[l]{\#},morestring=[b]"}
\lstdefinelanguage{sh}{morekeywords={cmake,ctest,cpack,find,sha256sum,clang-tidy,echo,cp,make,ninja,git},morecomment=[l]{\#},morestring=[b]"}
\lstdefinelanguage{yaml}{morecomment=[l]{\#},morestring=[b]",morekeywords={name,on,permissions,jobs,steps,uses,run,with,env,matrix,include,if}}
\lstdefinelanguage{text}{}
\lstset{basicstyle=\ttfamily\footnotesize,keywordstyle=\color{navy}\bfseries,commentstyle=\color{muted},stringstyle=\color{teal},backgroundcolor=\color{codegray},frame=single,rulecolor=\color{pale},breaklines=true,breakatwhitespace=false,showstringspaces=false,columns=fullflexible,keepspaces=true,tabsize=2,xleftmargin=5pt,xrightmargin=5pt,aboveskip=9pt,belowskip=9pt}
\providecommand{\tightlist}{\setlength{\itemsep}{0pt}\setlength{\parskip}{0pt}}
\providecommand{\passthrough}[1]{#1}
\providecommand{\real}[1]{#1}

\begin{document}
\hypersetup{pageanchor=false}\begin{titlepage}
\color{navy}
\vspace*{1.15in}
{\Large\bfseries PROFESSIONAL CMAKE\\[7pt]SOFTWARE BUILD ENGINEERING\par}
\vspace{0.32in}\noindent\textcolor{teal}{\rule{0.9\linewidth}{2.3pt}}\par
\vspace{0.45in}
{\parbox{0.95\textwidth}{\raggedright\fontsize{24}{29}\selectfont\bfseries CI/CD and DevSecOps Integration\par}}
\vspace{0.3in}
{\Large\color{teal}Build provenance, GitHub Actions, CodeQL\par}
\vfill
{\large DOCUMENT 06 OF 08\par}
\vspace{0.1in}
{\small Companion synthesis to Craig Scott, \emph{Professional CMake: A Practical Guide}, 20th Edition (2025).\\
Original learning/engineering material. Not affiliated with or endorsed by the author.\par}
\vspace{0.2in}
{\small 17 October 2026\par}
\end{titlepage}
\hypersetup{pageanchor=true}\clearpage\pagenumbering{roman}\tableofcontents\clearpage\pagenumbering{arabic}

\textbf{Purpose.} Connect CMake, CTest and CPack to a least-privilege
software delivery pipeline with measurable release controls.
\textbf{{[}BOOK-GROUNDED{]}} presets, package creation, static/dynamic
analysis and cross-platform builds are grounded in Chapters 24, 26--33,
35--37 and 42. \textbf{{[}ENGINEERING EXTENSION{]}} GitHub Actions
workflow policy, dependency governance, CodeQL, SBOM,
provenance/attestation, signing and organizational change controls are
extensions, not CMake features.

\chapter{Delivery pipeline and separation of
duties}\label{delivery-pipeline-and-separation-of-duties}

\begin{lstlisting}
      Pull Request        Protected Branch          Release Candidate
          |                      |                        |
          v                      v                        v
  lint / configure / test -> full OS/compiler matrix -> stage + consume
          |                      |                        |
      static analysis      dynamic analyzers         package & manifest
          |                      |                        |
      code scanning       dependency inventory      sign + attest + publish
          |                      |                        |
          +-------------- evidence + policy gates --------+
\end{lstlisting}

\textbf{Trust boundary:} source checkout, build scripts, package
managers, action dependencies, CMake providers, compilers and test
programs are executable inputs. Running them on a privileged release
runner grants code execution. Keep untrusted PR builds separate from
credentials and protected publishing environments.

\chapter{Pipeline stages and evidence
contracts}\label{pipeline-stages-and-evidence-contracts}

\begin{longtable}[]{@{}
  >{\raggedright\arraybackslash}p{(\columnwidth - 6\tabcolsep) * \real{0.2500}}
  >{\raggedright\arraybackslash}p{(\columnwidth - 6\tabcolsep) * \real{0.2500}}
  >{\raggedright\arraybackslash}p{(\columnwidth - 6\tabcolsep) * \real{0.2500}}
  >{\raggedright\arraybackslash}p{(\columnwidth - 6\tabcolsep) * \real{0.2500}}@{}}
\toprule\noalign{}
\begin{minipage}[b]{\linewidth}\raggedright
ID
\end{minipage} & \begin{minipage}[b]{\linewidth}\raggedright
Stage
\end{minipage} & \begin{minipage}[b]{\linewidth}\raggedright
Required output
\end{minipage} & \begin{minipage}[b]{\linewidth}\raggedright
Security constraint
\end{minipage} \\
\midrule\noalign{}
\endhead
\bottomrule\noalign{}
\endlastfoot
CI-01 & checkout and bootstrap & immutable commit SHA, versions &
pinned, reviewed action and toolchain provenance \\
CI-02 & configure & CMake cache/configuration logs & safe external
dependency policy; no secrets printed \\
CI-03 & compile/link & binaries and full compiler diagnostics &
consistent compiler, linker and sanitizer variants \\
CI-04 & execute CTest & machine-readable JUnit results & resource
isolation and non-prod credentials \\
CI-05 & static analysis & findings and coverage inventory & documented
analysis scope, no silent skip \\
CI-06 & stage + consumer & relocatable SDK and consumer result & no
source/build path leak \\
CI-07 & package & archive, manifest, checksums & deterministic input and
artifact review \\
CI-08 & attest/publish & provenance, signatures, registry coordinates &
protected environment; least-privilege OIDC \\
\end{longtable}

\chapter{Repository workflow example}\label{repository-workflow-example}

The included reference project contains
\passthrough{\lstinline!.github/workflows/build.yml!}. This standalone
illustration assumes Ubuntu, macOS and Windows runners have recent
CMake, Ninja and C++ compilers; real enterprise workflows pin versions
and validate that the runner images satisfy prerequisites.

\begin{lstlisting}
name: cmake-verify
on:
  push:
  pull_request:
permissions:
  contents: read
jobs:
  test:
    strategy:
      fail-fast: false
      matrix:
        os: [ubuntu-latest, windows-latest, macos-latest]
    runs-on: ${{ matrix.os }}
    steps:
      - uses: actions/checkout@v4
      - name: Configure
        run: cmake --preset ci-debug
      - name: Build
        run: cmake --build --preset ci-debug --parallel 2
      - name: Test
        run: ctest --preset ci-debug --output-on-failure
\end{lstlisting}

\textbf{Important:} runner labels and action tags are moving inputs. For
a hardened deployment, pin GitHub Actions to audited commit SHAs,
explicitly install/verify CMake/Ninja versions, and restrict
\passthrough{\lstinline!pull\_request\_target!} usage. This example is
instructional, not a security-hardened production workflow.

\chapter{Consistency between local and CI
builds}\label{consistency-between-local-and-ci-builds}

The same \passthrough{\lstinline!ci-debug!} or
\passthrough{\lstinline!ci-release!} preset should drive local
engineering and CI. Avoid YAML-specific compiler flags that diverge from
the preset. If CI requires additional settings, create a documented CI
variant rather than silently mutating defaults. Build caching can
accelerate development but must never be trusted as proof of source or
artifact integrity.

\begin{lstlisting}[language=sh]
cmake --preset ci-release
cmake --build --preset ci-release --parallel 4
ctest --preset ci-release --output-on-failure
cmake --install build/ci-release --prefix /tmp/buildsuite-install
\end{lstlisting}

\chapter{CodeQL C/C++ scanning
extension}\label{codeql-cc-scanning-extension}

As documented by GitHub, C/C++ supports \passthrough{\lstinline!none!},
\passthrough{\lstinline!autobuild!} and \passthrough{\lstinline!manual!}
build modes; exact availability and semantics depend on GitHub
product/version. If generated sources or custom configurations are
important, explicitly analyze whether \passthrough{\lstinline!none!} or
\passthrough{\lstinline!autobuild!} covers the relevant translation
units. A manual-build advanced-setup approach can invoke the same CMake
preset that CI uses.

\begin{lstlisting}
# Selected fragment for a separate code-scanning job:
permissions:
  contents: read
  security-events: write
steps:
  - uses: actions/checkout@v4
  - uses: github/codeql-action/init@v4
    with:
      languages: c-cpp
      build-mode: manual
  - run: cmake --preset ci-debug
  - run: cmake --build --preset ci-debug
  - uses: github/codeql-action/analyze@v4
\end{lstlisting}

\textbf{Governance:} Code scanning and build verification are distinct
controls. When policy requires scan-before-merge, ensure the required
analysis check belongs to the change's commit/PR head and the target's
branch protection rule. A scan after merging does not validate pre-merge
control effectiveness.

\chapter{Third-party dependency and source acquisition
controls}\label{third-party-dependency-and-source-acquisition-controls}

\begin{longtable}[]{@{}
  >{\raggedright\arraybackslash}p{(\columnwidth - 4\tabcolsep) * \real{0.3333}}
  >{\raggedright\arraybackslash}p{(\columnwidth - 4\tabcolsep) * \real{0.3333}}
  >{\raggedright\arraybackslash}p{(\columnwidth - 4\tabcolsep) * \real{0.3333}}@{}}
\toprule\noalign{}
\begin{minipage}[b]{\linewidth}\raggedright
Acquisition
\end{minipage} & \begin{minipage}[b]{\linewidth}\raggedright
Technical advantage
\end{minipage} & \begin{minipage}[b]{\linewidth}\raggedright
Policy requirement
\end{minipage} \\
\midrule\noalign{}
\endhead
\bottomrule\noalign{}
\endlastfoot
\passthrough{\lstinline!find\_package(CONFIG)!} & uses supplied package
binary and imported target metadata & approved package prefix, exact
version, provenance \\
\passthrough{\lstinline!FetchContent!} & integrates dependency targets
in normal build & immutable source revision and code review; network
restrictions \\
\passthrough{\lstinline!ExternalProject!} & isolates independent
configure/build workflows & artifact handoff, checksum and ownership \\
vcpkg/Conan/package manager & centralized resolution / binary caching &
version lockfiles, mirrors, verified metadata \\
vendored source & fixed repository revision & licensing, update
ownership, recursive scan \\
\end{longtable}

\textbf{Control requirement:} capture an inventory including name,
version or commit, source location, license, declared direct/transitive
status, runtime distribution relevance and vulnerability advisory
mapping. Produce an SPDX or CycloneDX SBOM with an appropriate
\textbf{external} tool and verify its completeness; CMake itself does
not certify an SBOM.

\section{Dependency exception record}\label{dependency-exception-record}

\begin{longtable}[]{@{}
  >{\raggedright\arraybackslash}p{(\columnwidth - 2\tabcolsep) * \real{0.5000}}
  >{\raggedright\arraybackslash}p{(\columnwidth - 2\tabcolsep) * \real{0.5000}}@{}}
\toprule\noalign{}
\begin{minipage}[b]{\linewidth}\raggedright
Field
\end{minipage} & \begin{minipage}[b]{\linewidth}\raggedright
Evidence to record
\end{minipage} \\
\midrule\noalign{}
\endhead
\bottomrule\noalign{}
\endlastfoot
Exception ID and owner & ticket, accountable engineering team \\
Dependency coordinate & repository / package URL and immutable
version \\
Why necessary & functionality, competing alternatives, build impact \\
Security posture & advisories, active exploitation if known,
compensating controls \\
License & SPDX expression and distribution obligations \\
Approval expiry & trigger for re-evaluation and planned removal \\
\end{longtable}

\chapter{Software bill of materials, provenance and artifact
integrity}\label{software-bill-of-materials-provenance-and-artifact-integrity}

\textbf{{[}ENGINEERING EXTENSION{]}} CMake can generate install
manifests and archives but does not establish cryptographic artifact
origin by itself. An enterprise release pipeline should collect: source
revision, build workflow identity, resolved dependencies,
compiler/toolchain metadata, tests and scan results, artifact hashes and
policy decisions. Sign final artifacts or issue attestations with
organization-managed identity; verify before promotion.

Threat coverage:

\begin{itemize}
\tightlist
\item
  \emph{Dependency confusion / package substitution:} constrain sources
  and verify immutable artifacts.
\item
  \emph{Compromised CI action / script:} pin reviewable versions, limit
  permissions, separate privileged jobs.
\item
  \emph{Toolchain tampering:} use audited toolchain images and record
  provenance.
\item
  \emph{Secret exposure:} inject only into protected publish jobs, not
  configure/test of untrusted PRs.
\item
  \emph{Artifact substitution:} hash, attest/sign and verify registry
  artifacts at deployment.
\item
  \emph{Malicious test fixtures:} sandbox tests and restrict outbound
  network and filesystem access.
\end{itemize}

\chapter{SLSA-style capability
roadmap}\label{slsa-style-capability-roadmap}

The following maturity model is \textbf{inspired by} supply-chain
assurance practice, not a claim of formal SLSA compliance:

\begin{longtable}[]{@{}
  >{\raggedright\arraybackslash}p{(\columnwidth - 4\tabcolsep) * \real{0.3333}}
  >{\raggedright\arraybackslash}p{(\columnwidth - 4\tabcolsep) * \real{0.3333}}
  >{\raggedright\arraybackslash}p{(\columnwidth - 4\tabcolsep) * \real{0.3333}}@{}}
\toprule\noalign{}
\begin{minipage}[b]{\linewidth}\raggedright
Level
\end{minipage} & \begin{minipage}[b]{\linewidth}\raggedright
Build characteristics
\end{minipage} & \begin{minipage}[b]{\linewidth}\raggedright
Promotion criterion
\end{minipage} \\
\midrule\noalign{}
\endhead
\bottomrule\noalign{}
\endlastfoot
0: ad hoc & local developer binaries & no unattended release \\
1: repeatable & declared presets, automated build/test &
machine-readable evidence \\
2: controlled & protected CI, versioned toolchains, dependency inventory
& risk-based gates \\
3: verifiable & isolated builder, signed provenance and manifest &
verification in staging \\
4: defended & policy-controlled dependencies, continuous validation,
exception management & auditable change and release records \\
\end{longtable}

Formal compliance with a specific SLSA version requires checking its
published requirements separately.

\chapter{Release manifest and reproducibility
measurement}\label{release-manifest-and-reproducibility-measurement}

Record each build's configuration and environment without claiming
bit-for-bit identical results until tested. Sources of nondeterminism
include timestamps, filesystem paths, locale, archive ordering,
randomness, incremental object reuse and toolchain versions. A
reproducibility experiment should produce two independent clean builds
from the same inputs, compare artifact hashes and investigate
differences using structured binary-diff tooling.

\begin{lstlisting}[language=sh]
sha256sum path/to/release-artifact.tar.gz
# Repeat after a second isolated clean build.
\end{lstlisting}

Hash equality in one scenario does not prove resilience against a
malicious compiler or compromised package repository.

\chapter{Release promotion checklist}\label{release-promotion-checklist}

\begin{itemize}
\tightlist
\item[$\square$]
  Branch protection requires all declared build, test and security
  statuses.
\item[$\square$]
  CI uses declared presets, compiler versions and separate Debug/Release
  builds.
\item[$\square$]
  Required static and dynamic analysis jobs did not silently skip
  checks.
\item[$\square$]
  CodeQL mode and coverage are documented for compiled languages.
\item[$\square$]
  Dependencies are pinned, inventoried, licensed and reviewed.
\item[$\square$]
  A clean installed consumer validates relocatability and compatibility.
\item[$\square$]
  Packages are scanned; manifests and hashes have been retained.
\item[$\square$]
  Publish permissions are isolated from untrusted builds and PRs.
\item[$\square$]
  Final artifacts are signed/attested as required by organizational
  policy.
\item[$\square$]
  Rollback version and supply-chain incident response owner are known.
\end{itemize}

\chapter{Incident / failing build
response}\label{incident-failing-build-response}

\begin{enumerate}
\def\labelenumi{\arabic{enumi}.}
\tightlist
\item
  Preserve failing run IDs, logs, build matrix and commit state.
\item
  Quarantine any potentially compromised artifact; suspend promotion.
\item
  Compare dependency versions and toolchain identity against approved
  baseline.
\item
  Re-run in an isolated runner with controlled network access.
\item
  Determine whether incident is configure-time execution, dependency
  drift, faulty source, infrastructure or policy gap.
\item
  Record causal finding, rebuild trusted artifacts and reconcile
  SBOM/provenance.
\item
  Restore release flow only after independent validation and documented
  approval.
\end{enumerate}

\chapter{Official references}\label{official-references}

CMake:
\url{https://cmake.org/cmake/help/latest/manual/cmake-presets.7.html} ;
\url{https://cmake.org/cmake/help/latest/manual/cpack.1.html} . GitHub
CodeQL:
\url{https://docs.github.com/en/code-security/concepts/code-scanning/codeql/codeql-for-compiled-languages}
. GitHub Actions: \url{https://docs.github.com/en/actions} . SPDX:
\url{https://spdx.dev/} ; CycloneDX: \url{https://cyclonedx.org/} ;
SLSA: \url{https://slsa.dev/} . Revalidate versions and exact
capabilities when implementing a real production pipeline.

\end{document}
